\begin{pastexam}[label={px:ncku-112-prog-07}]{成大資訊聯招 112 學年度 第 7 題}
（選擇）

{[5\%]} A hash table in size of 10 adopts open addressing with linear probing, where its hash function
is defined as h(k)=k\%10. The table is initially empty. We insert 6 data items into the table, resulting
in the following table. Which of the following choices gives a possible order for inserting the data
items to the table?
\begin{center}
\begin{tabular}{|c|c|}
\hline
0 & \\ \hline
1 & \\ \hline
2 & 22 \\ \hline
3 & 43 \\ \hline
4 & 14 \\ \hline
5 & 32 \\ \hline
6 & 26 \\ \hline
7 & 53 \\ \hline
8 & \\ \hline
9 & \\ \hline
\end{tabular}
\end{center}
(A)~32, 14, 53, 26, 22, 43\\
(B)~42, 22, 13, 26, 53, 32\\
(C)~22, 43, 14, 32, 53, 26\\
(D)~43, 22, 14, 26, 32, 53\\
(E)~14, 43, 22, 32, 26, 53
\end{pastexam}

\begin{pastsoln}
答案：(D)(E)。

線性探測（資結書第~12.3~節）從 $h(k)$ 開始往後找第一個空格。 一個鍵落在 $h(k) + t$，表示它插入時
$h(k), \ldots, h(k) + t - 1$ 都已經有鍵，而 $h(k) + t$ 是空的。 逐個看：
\begin{center}
\small
\begin{tabular}{cccl}
\toprule
鍵 & $h(k)$ & 所在格 & 插入前必須已經有鍵的格 \\
\midrule
$22$ & $2$ & $2$ & 無 \\
$43$ & $3$ & $3$ & 無 \\
$14$ & $4$ & $4$ & 無 \\
$32$ & $2$ & $5$ & $2, 3, 4$（$22, 43, 14$） \\
$26$ & $6$ & $6$ & 無 \\
$53$ & $3$ & $7$ & $3, 4, 5, 6$（$43, 14, 32, 26$） \\
\bottomrule
\end{tabular}
\end{center}
所以 $32$ 要在 $22, 43, 14$ 之後，$53$ 要在 $43, 14, 32, 26$ 之後。 另外，落在自己家（$h(k)$）的鍵，插入時那一格
必須是空的；會搶走這些格子的只有 $32$（從 $2$ 開始找）和 $53$（從 $3$ 開始找），而它們都排在後面，所以上面
兩個條件就夠了。

\begin{itemize}
\item (A) $32$ 第一個插入，會直接放在第 $2$ 格。 錯。
\item (B) 有 $42$、$13$，表裡沒有這兩個數（而且少了 $43$、$14$）。 錯。
\item (C) $53$ 在 $26$ 之前：插入 $53$ 時第 $3$--$5$ 格有鍵、第 $6$ 格是空的，$53$ 放進第 $6$ 格，之後 $26$ 只能
放第 $7$ 格，和表不同。 錯。
\item (D) $43, 22, 14, 26, 32, 53$：$32$ 在 $22, 43, 14$ 之後；$53$ 在最後。 照著插一遍得到的表和題目一樣。 對。
\item (E) $14, 43, 22, 32, 26, 53$：$32$ 在前三個之後，$53$ 在最後。 對。
\end{itemize}
\end{pastsoln}
